\subsection{Bonus DTL}
\label{subsec:dtl-bonus}

\todo{Hapus subbab bonus yang tidak dikerjakan.}

\subsubsection{Visualisasi Percabangan Pohon}
\todo{Gambar percabangan pohon beserta parameter \textit{top-N} percabangan yang ditampilkan. Jelaskan fungsi visualisasi, cara memilih \textit{top-N}, dan tampilkan gambarnya.}

\begin{figure}[H]
	\centering
	\placeholdergambar{Visualisasi Percabangan Pohon}
	\caption{Visualisasi Percabangan Pohon DTL}
	\label{fig:dtl-bonus-0}
\end{figure}

\subsubsection{Algoritma Optimasi di Luar Materi Kelas}
\todo{Jika ada: jelaskan algoritmanya, mengapa dapat mengoptimasi model, bukti perbaikannya (sebelum vs sesudah), dan \textbf{referensinya (tanpa referensi bonus dianulir)}. Misalnya \textit{cost-complexity pruning}, penalti MDL pada \textit{threshold} kontinu, \textit{shrinkage} probabilitas daun, atau Laplace \textit{correction}.}
